\scanpage{65}
\[
H^*(X_{\mathrm{cris}},F)\simeq H^*(X_{\mathrm{zar}},\mathcal F^*)
\]
d'où
\[
H^*(X_{\mathrm{cris}},F)\Longleftarrow H^p(\nu\longmapsto H^q(X,\mathcal F^{(\nu)}))
\]
6.3. \uline{Cas particulier} Supposons que \deleted{\unsure{$X\subset Y$ est}}\inserted{l'on ait $X_0\subset X$} une immersion \uline{nilpotente}, \deleted{\unsure{Alors}} et que $F$ sur $X_{\mathrm{cris}}$ provienne d'un $F_0$ sur $X_{0\,\mathrm{cris}}$ par ``restriction''. Alors $\forall$ entier $\nu$, les complétés formels $Y^{\nu+1}_{/X_0}$ et $Y^{\nu+1}_{/X}$ sont identiques, On trouve donc
\margin{\uline{Application aux}\par\uline{\unsure{stratifications}}\par\uline{canoniques}}
\[
\mathcal F^{(\nu)}\simeq\mathcal F_0^{(\nu)}\quad\uline{\unsure{can.}}\quad\text{sur }X_{\mathrm{zar}}=X_{0\,\mathrm{zar}},
\]
d'où aussitôt
\[
H^*(X_{\mathrm{cris}},F)\xrightarrow{\sim}H^*(X_{0\,\mathrm{cris}},F_0).
\]
Ceci a été établi lorsqu'on peut plonger $X$ globalement dans un $Y$ \unsure{f\textsuperscript{t}} lisse sur $S$, mais \deleted{\unsure{par un rais. local}} utilisant un rec. de $X$ par des ouverts affines, on voit que l'on peut se débarrasser de cette hypothèse.

[Il faudrait regarder si on peut se débarrasser de l'hypothèse type quasi-cohérence faite sur $F$ ....?]
\mspage{64}
